% ============ 第1章 §1.10 Hilbert变换的L²有界性(再访)(原书 pp.78-88) ============
\section{\texorpdfstring{Hilbert 变换的 $L^2$ 有界性(再访)}{Hilbert 变换的 L2 有界性(再访)}}\label{sec:1.10}
\index{Hilbert 变换!L2 有界性}


Hilbert 变换的 $L^{2}$ 有界性是用 Fourier 变换证明的,但在满足与 Hilbert 变换类似的积分核条件但无齐次性的情形下,有时不能使用 Fourier 变换.这种情形下,作为证明 $L^{2}$ 有界性的有效方法,人们开发了称为 Cotlar--Knapp--Stein 概正交性引理的方法.本节介绍它,并举例说明用它可证明 Hilbert 变换的 $L^{2}$ 有界性.

\begin{lemma}[Cotlar--Knapp--Stein 概正交性引理]\label{lem:1.10.1}
设 $\{T_{j}\}_{j \in \mathbb{Z}}$ 是 Hilbert 空间 $H$ 上的有界线性算子列,设存在数列 $\{a(j)\}_{j \in \mathbb{Z}}$ 满足
\begin{equation}\tag{1.10.1}
\| T_{j}^{*}T_{k} \|_{L(H)} + \| T_{j}T_{k}^{*} \|_{L(H)} \leq a(j - k) \quad (j, k \in \mathbb{Z})
\label{eq:1.10.1}\end{equation}
与
\begin{equation}\tag{1.10.2}
A := \sum_{j \in \mathbb{Z}}\sqrt{a(j)} < \infty.
\label{eq:1.10.2}\end{equation}
则下面结论成立:

(i) 对 $\mathbb{Z}$ 的任意有限子集 $\Lambda$,
\begin{equation*}
\left\| \sum_{j \in \Lambda}T_{j} \right\|_{L(H)} \leq A.
\end{equation*}

(ii) 对所有 $x \in H$,
\begin{equation*}
\sum_{j \in \mathbb{Z}}\| T_{j}(x) \|_{H}^{2} \leq A^{2}\| x \|_{H}^{2}.
\end{equation*}

(iii) 对每个 $x \in H$,当 $N \to \infty$ 时 $\sum_{|j| \leq N}T_{j}(x)$ 依 $H$ 范数收敛到某个 $T(x) \in H$,如此定义的 $T$ 是 $H$ 上的有界线性算子,且
\begin{equation*}
\| T(x) \|_{H} \leq A\| x \|_{H} \quad (x \in H).
\end{equation*}

其中 $T_{j}^{*}$ 是 $T_{j}$ 的伴随算子,$\| T_{j} \|_{L(H)}$ 表示 $T_{j}$ 的算子范数.
\end{lemma}

\begin{proof}
(i) 一般地,对 $H$ 上的有界线性算子 $S$,$\| S \|_{L(H)}^{2} = \| SS^{*} \|_{L(H)}$,故 $\| S \|_{L(H)}^{2} = \| [SS^{*}]^{2^{\ell}} \|_{L(H)}^{2^{-\ell}}$($\ell \in \mathbb{N}$),于是
\begin{equation*}
B^{2} := \left\| \sum_{j \in \Lambda}T_{j} \right\|_{L(H)}^{2} = \left\| \left[\left(\sum_{j \in \Lambda}T_{j}\right)\left(\sum_{j \in \Lambda}T_{j}^{*}\right)\right]^{2^{\ell}} \right\|_{L(H)}^{2^{-\ell}}.
\end{equation*}
因此令 $m = 2^{\ell}$,展开右边范数记号内的部分,
\begin{equation}\tag{1.10.3}\begin{aligned}
B^{2} &= \left\| \displaystyle\sum_{j_{1}, j_{2}, \dots, j_{2m} \in \Lambda}T_{j_{1}}T_{j_{2}}^{*}\cdots T_{j_{2m-1}}T_{j_{2m}}^{*} \right\|_{L(H)}^{1/m} \\
&\leq \left(\displaystyle\sum_{j_{1}, j_{2}, \dots, j_{2m} \in \Lambda}\left\| T_{j_{1}}T_{j_{2}}^{*}\cdots T_{j_{2m-1}}T_{j_{2m}}^{*} \right\|_{L(H)}\right)^{1/m}.
\end{aligned}\label{eq:1.10.3}\end{equation}
由 \eqref{eq:1.10.1},
\begin{align*}
\left\| T_{j_{1}}T_{j_{2}}^{*}\cdots T_{j_{2m-1}}T_{j_{2m}}^{*} \right\|_{L(H)} & \leq \| T_{j_{1}}T_{j_{2}}^{*} \|_{L(H)}\cdots\| T_{j_{2m-1}}T_{j_{2m}}^{*} \|_{L(H)} \\
& \leq a(j_{1} - j_{2})a(j_{3} - j_{4})\cdots a(j_{2m-1} - j_{2m}),
\end{align*}
且
\begin{align*}
& \left\| T_{j_{1}}T_{j_{2}}^{*}\cdots T_{j_{2m-1}}T_{j_{2m}}^{*} \right\|_{L(H)} \\
& \quad \leq \| T_{j_{1}} \|_{L(H)}\| T_{j_{2}}^{*}T_{j_{3}} \|_{L(H)}\cdots\| T_{j_{2m-2}}^{*}T_{j_{2m-1}} \|_{L(H)}\| T_{j_{2m}}^{*} \|_{L(H)} \\
& \quad \leq \sqrt{a(0)}a(j_{2} - j_{3})a(j_{4} - j_{5})\cdots a(j_{2m-2} - j_{2m-1})\sqrt{a(0)},
\end{align*}
故取两者的几何平均,
\begin{align*}
\left\| T_{j_{1}}T_{j_{2}}^{*}\cdots T_{j_{2m-1}}T_{j_{2m}}^{*} \right\|_{L(H)}&\leq \sqrt{a(0)}\sqrt{a(j_{1} - j_{2})}\sqrt{a(j_{2} - j_{3})}\cdots\sqrt{a(j_{2m-1} - j_{2m})}.
 \end{align*}
于是由 \eqref{eq:1.10.2},
\begin{align*}
& \displaystyle\sum_{j_{1}, j_{2}, \ldots, j_{2m} \in \Lambda}\left\| T_{j_{1}}T_{j_{2}}^{*}\cdots T_{j_{2m-1}}T_{j_{2m}}^{*} \right\|_{L(H)} \\
& \leq \sqrt{a(0)}\displaystyle\sum_{j_{2}, \ldots, j_{2m} \in \Lambda}\left(\sum_{j_{1} \in \Lambda}\sqrt{a(j_{1} - j_{2})}\right)\sqrt{a(j_{2} - j_{3})}\cdots\sqrt{a(j_{2m-1} - j_{2m})} \\
& \ldots \\
& \leq \sqrt{a(0)}A^{2m-1}\sum_{j_{2m} \in \Lambda}1 = \sqrt{a(0)}A^{2m-1}(\Lambda \text{ 的元素个数}).
\end{align*}
与 \eqref{eq:1.10.3} 合并,
\begin{equation*}
B^{2} \leq a(0)^{1/(2m)}A^{2-1/m}(\Lambda \text{ 的元素个数})^{1/m}.
\end{equation*}
令 $\ell \to \infty$,则 $m = 2^{\ell} \to \infty$,故
\begin{equation*}
\left\| \sum_{j \in \Lambda}T_{j} \right\|_{L(H)} = B \leq A.
\end{equation*}

(ii) 设 $x \in H$.对实数 $\theta$,算子列 $\{e^{2\pi ij\theta}T_{j}\}_{j \in \mathbb{Z}}$ 也满足 \eqref{eq:1.10.1},\eqref{eq:1.10.2},故由 (i),
\begin{equation}\tag{1.10.4}
\left\| \sum_{|j| \leq N}e^{2\pi ij\theta}T_{j}(x) \right\|_{H} \leq A\| x \|_{H}.
\label{eq:1.10.4}\end{equation}
另一方面,
\begin{align*}
\displaystyle\int_{0}^{1}\left\| \sum_{|j| \leq N}e^{2\pi ij\theta}T_{j}(x) \right\|_{H}^{2}\dd\theta&= \displaystyle\int_{0}^{1}\left(\sum_{|j| \leq N}e^{2\pi ij\theta}T_{j}(x), \sum_{|k| \leq N}e^{2\pi ik\theta}T_{k}(x)\right)\dd\theta  \\
&= \displaystyle\sum_{|j| \leq N}\sum_{|k| \leq N}\int_{0}^{1}e^{2\pi i(j-k)\theta}\dd\theta\left(T_{j}(x), T_{k}(x)\right) \\
&= \displaystyle\sum_{|j| \leq N}\| T_{j}(x) \|_{H}^{2}.
\end{align*}
与 \eqref{eq:1.10.4} 合并,
\begin{equation*}
\sum_{|j| \leq N}\| T_{j}(x) \|_{H}^{2} \leq A^{2}\| x \|_{H}^{2}. \quad \text{令 } N \to \infty \text{ 得 } \sum_{j \in \mathbb{Z}}\| T_{j}(x) \|_{H}^{2} \leq A^{2}\| x \|_{H}^{2}.
\end{equation*}

(iii) 设 $x \in H$.首先证明 $\left\{\sum_{|j| \leq N}T_{j}(x)\right\}_{N=1}^{\infty}$ 是 $H$ 中的 Cauchy 列.若否定这一点,则存在正数 $\varepsilon$ 与自然数列 $N_{1} < N_{2} < \cdots$ 使
\begin{equation}\tag{1.10.5}
\left\| \sum_{N_{k} < |j| \leq N_{k+1}}T_{j}(x) \right\|_{H} \geq \varepsilon \quad (k = 1, 2, \dots).
\label{eq:1.10.5}\end{equation}
另一方面,令
\begin{equation*}
S_{j} = e^{2\pi ik\theta}T_{j} \quad (N_{k} < |j| \leq N_{k+1},\ k = 1, 2, \dots),
\end{equation*}
容易看出 $\{S_{j}\}_{j \in \mathbb{Z}}$ 也满足 \eqref{eq:1.10.1},\eqref{eq:1.10.2},故由 (i),
\begin{equation}\tag{1.10.6}
\left\| \sum_{k=1}^{K}e^{2\pi ik\theta}\sum_{N_{k} < |j| \leq N_{k+1}}T_{j}(x) \right\|_{H} \leq A\| x \|_{H}.
\label{eq:1.10.6}\end{equation}
与 (ii) 中的计算一样,
\begin{equation}\tag{1.10.7}
\int_{0}^{1}\left\| \sum_{k=1}^{K}e^{2\pi ik\theta}\sum_{N_{k} < |j| \leq N_{k+1}}T_{j}(x) \right\|_{H}^{2}\dd\theta = \sum_{k=1}^{K}\left\| \sum_{N_{k} < |j| \leq N_{k+1}}T_{j}(x) \right\|_{H}^{2}.
\label{eq:1.10.7}\end{equation}
由 \eqref{eq:1.10.5},\eqref{eq:1.10.6},\eqref{eq:1.10.7},
\begin{equation*}
K\varepsilon^{2} \leq \sum_{k=1}^{K}\left\| \sum_{N_{k} < |j| \leq N_{k+1}}T_{j}(x) \right\|_{H}^{2} \leq A^{2}\| x \|_{H}^{2},
\end{equation*}
令 $K \to \infty$ 即得矛盾.因此 $\left\{\sum_{|j| \leq N}T_{j}(x)\right\}_{N=1}^{\infty}$ 是 $H$ 中的 Cauchy 列.由 Hilbert 空间的完备性,当 $N \to \infty$ 时 $\sum_{|j| \leq N}T_{j}(x)$ 依 $H$ 范数收敛到某个 $T(x) \in H$.另外,对 $\alpha, \beta \in \mathbb{C}$,$x, y \in H$,
\begin{align*}
\left\| \displaystyle\sum_{|j| \leq N}T_{j}(\alpha x + \beta y) - \alpha T(x) - \beta T(y) \right\|_{H}&= \left\| \alpha\left(\displaystyle\sum_{|j| \leq N}T_{j}(x) - T(x)\right) + \beta\left(\sum_{|j| \leq N}T_{j}(y) - T(y)\right) \right\|_{H}  \\
&\leq |\alpha|\left\| \displaystyle\sum_{|j| \leq N}T_{j}(x) - T(x) \right\|_{H} + |\beta|\left\| \sum_{|j| \leq N}T_{j}(y) - T(y) \right\|_{H},
\end{align*}
故当 $N \to \infty$ 时 $\sum_{|j| \leq N}T_{j}(\alpha x + \beta y)$ 收敛到 $\alpha T(x) + \beta T(y)$.另一方面,$\sum_{|j| \leq N}T_{j}(\alpha x + \beta y)$ 本来收敛到 $T(\alpha x + \beta y)$,故 $T(\alpha x + \beta y) = \alpha T(x) + \beta T(y)$,即 $T$ 是从 $H$ 到 $H$ 的线性算子.进而,
\begin{align*}
\| T(x) \|_{H} & \leq \left\| T(x) - \displaystyle\sum_{|j| \leq N}T_{j}(x) \right\|_{H} + \left\| \sum_{|j| \leq N}T_{j}(x) \right\|_{H} \\
& \leq \left\| T(x) - \displaystyle\sum_{|j| \leq N}T_{j}(x) \right\|_{H} + A\| x \|_{H},
\end{align*}
故令 $N \to \infty$ 得
\begin{equation*}
\| T(x) \|_{H} \leq A\| x \|_{H}.
\end{equation*}
因此 $T$ 是 $H$ 上的有界线性算子.
\end{proof}

利用这个引理,可证明如下奇异积分的 $L^{2}$ 有界性.

\begin{theorem}\label{thm:1.10.2}
对二元函数 $K(x, y)$,设存在满足以下 4 个条件的正数 $C_{1}, C_{2}, C_{3}$ 与 $0 < \delta \leq 1$:

(i)
\begin{equation*}
|K(x, y)| \leq \frac{C_{1}}{|x - y|} \quad (x \neq y),
\end{equation*}

(ii)
\begin{equation*}
|K(x, y) - K(x, z)| \leq \frac{C_{2}}{|x - y|}\left(\frac{|y - z|}{|x - y|}\right)^{\delta} \quad (|y - z| < |x - y|/2),
\end{equation*}
\begin{equation*}
|K(x, y) - K(z, y)| \leq \frac{C_{3}}{|x - y|}\left(\frac{|x - z|}{|x - y|}\right)^{\delta} \quad (|x - z| < |x - y|/2),
\end{equation*}

(iii)
\begin{equation*}
\lim_{\varepsilon \to 0}\int_{\varepsilon < |x - y| < 2^{m}}K(x, y)\dd y = 0 \quad (m \in \mathbb{Z},\ x \in \mathbb{R}),
\end{equation*}

(iv)
\begin{equation*}
\int_{2^{m} \leq |x - y| < 2^{m+1}}K(x, y)\dd x = 0 \quad (m \in \mathbb{Z},\ y \in \mathbb{R}).
\end{equation*}

则当 $f \in C^{1}(\mathbb{R})$ 且 $|f(x)| \leq C_{4}/(1 + |x|)^{2}$,$|f'(x)| \leq C_{5}/(1 + |x|)$ 时,
\begin{equation*}
Tf(x) := \mathrm{p.v.}\int K(x, y)f(y)\dd y = \lim_{\varepsilon \to 0}\int_{|x - y| > \varepsilon}K(x, y)f(y)\dd y
\end{equation*}
对每个 $x \in \mathbb{R}$ 存在,并对某个常数 $C_{6}$ 有
\begin{equation*}
|Tf(x)| \leq \frac{C_{6}}{1 + |x|},
\end{equation*}
且存在不依赖于 $f$ 的常数 $C$ 使
\begin{equation*}
\| Tf \|_{L^{2}} \leq C\| f \|_{L^{2}} \quad (f \in \mathcal{S}(\mathbb{R})).
\end{equation*}
\end{theorem}

\begin{proof}
注意由条件 (iii) 得
\begin{equation}\tag{1.10.8}
\int_{2^{m} \leq |x - y| < 2^{m+1}}K(x, y)\dd y = 0 \quad (m \in \mathbb{Z},\ x \in \mathbb{R}).
\label{eq:1.10.8}\end{equation}
现在,设 $f \in C^{1}(\mathbb{R})$ 且 $|f(x)| \leq C_{4}/(1 + |x|)^{2}$,$|f'(x)| \leq C_{5}/(1 + |x|)$,则
\begin{align*}
T_{\varepsilon}f(x)&:= \displaystyle\int_{|x - y| > \varepsilon}K(x, y)f(y)\dd y  \\
&= f(x)\displaystyle\int_{\varepsilon < |x - y| < 1}K(x, y)\dd y + \int_{\varepsilon < |x - y| < 1}K(x, y)\left(f(y) - f(x)\right)\dd y \\
&\quad + \displaystyle\int_{|x - y| \geq 1}K(x, y)f(y)\dd y.
\end{align*}
当 $0 < |x - y| \leq 1$ 时,
\begin{equation*}
|K(x, y)\left(f(y) - f(x)\right)| \leq C_{1}\max_{|x - z| \leq 1}|f'(z)| \leq \max_{|x - z| \leq 1}\frac{C_{1}C_{5}}{1 + |z|} \leq \frac{2C_{1}C_{5}}{1 + |x|}.
\end{equation*}
当 $|x - y| > 1$ 时,
\begin{equation*}
|K(x, y)f(y)| \leq \frac{C_{1}C_{4}}{|x - y|(1 + |y|)^{2}} \leq \frac{C_{1}C_{4}}{(1 + |y|)^{2}}.
\end{equation*}
故 $K(x, y)\left(f(y) - f(x)\right)$ 在 $[x - 1, x + 1]$ 上可积,$K(x, y)f(y)$ 在 $(-\infty, x - 1)$,$(x + 1, \infty)$ 上可积,故利用条件 (iii),
\begin{align*}
\lim_{\varepsilon \to 0}T_{\varepsilon}f(x) &= \displaystyle\int_{|x - y| < 1}K(x, y)\left(f(y) - f(x)\right)\dd y + \int_{|x - y| \geq 1}K(x, y)f(y)\dd y \\
&=: I_{1}(x) + I_{2}(x).
\end{align*}
以下与 Hilbert 变换的引理~\ref{lem:1.3.2} 一样,可证明存在常数 $C_{6}$ 使
\begin{equation*}
I_{1}(x) + I_{2}(x) \leq \frac{C_{6}}{1 + |x|} \quad (x \in \mathbb{R}).
\end{equation*}

这里,对每个 $j \in \mathbb{Z}$ 令
\begin{equation*}
K_{j}(x, y) := K(x, y)\chi_{\{2^{j} \leq |x - y| < 2^{j+1}\}}(x, y).
\end{equation*}
对每个 $j \in \mathbb{Z}$,$K_{j}(x, y)$ 作为 $y$ 的函数属于 $L^{2}(\mathbb{R})$,故当 $f \in L^{2}(\mathbb{R})$ 时可定义
\begin{equation}\tag{1.10.9}
T_{j}f(x) := \int_{\mathbb{R}}K_{j}(x, y)f(y)\dd y.
\label{eq:1.10.9}\end{equation}
对这个 $T_{j}$,
\begin{equation}\tag{1.10.10}
\int_{\mathbb{R}}|K_{j}(x, y)|\dd y,\ \int_{\mathbb{R}}|K_{j}(x, y)|\dd x \leq \int_{2^{j} \leq |x - y| < 2^{j+1}}\frac{C_{1}}{|x - y|}\dd y \leq 2C_{1}\log 2,
\label{eq:1.10.10}\end{equation}
故使用 Cauchy--Schwarz 不等式与 Fubini 定理,
\begin{align*}
\| T_{j}f \|_{L^{2}}&\leq \left(\displaystyle\int_{\mathbb{R}}\left[\left(\int_{\mathbb{R}}|K_{j}(x, y)|\dd y\right)^{1/2}\left(\int_{\mathbb{R}}|K_{j}(x, y)||f(y)|^{2}\dd y\right)^{1/2}\right]^{2}\dd x\right)^{1/2}  \\
&\leq \sqrt{2C_{1}\log 2}\left(\displaystyle\int_{\mathbb{R}}\int_{\mathbb{R}}|K_{j}(x, y)||f(y)|^{2}\dd y\dd x\right)^{1/2} \\
&\leq 2C_{1}\log 2\| f \|_{L^{2}}.
\end{align*}
因此
\begin{equation}\tag{1.10.11}
\| T_{j}T_{k}^{*} \|_{L(H)} + \| T_{j}^{*}T_{k} \|_{L(H)} \leq 8C_{1}^{2}\log^{2}2 \quad (j, k \in \mathbb{Z}).
\label{eq:1.10.11}\end{equation}
其次,对 $j, k \in \mathbb{Z}$ 令
\begin{equation*}
S_{j,k}(x, y) := \int_{\mathbb{R}}K_{j}(x, z)\overline{K_{k}(y, z)}\dd z,
\end{equation*}
则对 $f \in \mathcal{S}(\mathbb{R})$,
\begin{equation*}
T_{j}T_{k}^{*}f(x) = \int_{\mathbb{R}}K_{j}(x, z)\left(\int_{\mathbb{R}}\overline{K_{k}(y, z)}f(y)\dd y\right)\dd z = \int_{\mathbb{R}}S_{j,k}(x, y)f(y)\dd y.
\end{equation*}
于是来考察 $S_{j,k}(x, y)$.

(i) $k \leq j - 3$ 的情形.

(a) 当 $|x - y| \geq 2^{j+1} + 2^{k+1}$ 时,若 $2^{j} \leq |x - z| < 2^{j+1}$ 且 $2^{k} \leq |y - z| < 2^{k+1}$,则 $|x - y| \leq |x - z| + |y - z| < 2^{j+1} + 2^{k+1}$,故 $K_{j}(x, z)\overline{K_{k}(y, z)} = 0$,
\begin{equation*}
S_{j,k}(x, y) = 0.
\end{equation*}

(b) 当 $2^{j+1} - 2^{k+1} \leq |x - y| < 2^{j+1} + 2^{k+1}$ 时,若 $2^{k} \leq |y - z| < 2^{k+1}$,则 $|x - z| \geq |x - y| - |y - z| > 2^{j+1} - 2^{k+1} - 2^{k+1} > 2^{j}$,故
\begin{equation*}
|S_{j,k}(x, y)| \leq \int_{2^{k} \leq |y - z| < 2^{k+1}}\frac{C_{1}}{|x - z|}\frac{C_{1}}{|y - z|}\dd z \leq 2C_{1}^{2}2^{-j}\log 2.
\end{equation*}

(c) 当 $2^{j} + 2^{k+1} \leq |x - y| < 2^{j+1} - 2^{k+1}$ 时,若 $2^{k} \leq |y - z| < 2^{k+1}$,则 $|y - z| < |x - y|/2$,且 $|x - z| \geq |x - y| - |y - z| \geq 2^{j} + 2^{k+1} - 2^{k+1} = 2^{j}$,$|x - z| \leq |x - y| + |y - z| \leq 2^{j+1} - 2^{k+1} + 2^{k+1} = 2^{j+1}$,故 $K_{j}(x, z) = K(x, z)$.因此由 \eqref{eq:1.10.8},
\begin{equation*}
S_{j,k}(x, y) = \int_{\mathbb{R}}\left(K(x, z) - K(x, y)\right)\overline{K_{k}(y, z)}\dd z.
\end{equation*}
又已知 $|y - z| < |x - y|/2$,故由条件 (ii),
\begin{align*}
|S_{j,k}(x, y)| &\leq \displaystyle\int_{2^{k} \leq |y - z| < 2^{k+1}}\frac{C_{2}}{|x - y|}\left(\frac{|y - z|}{|x - y|}\right)^{\delta}\frac{C_{1}}{|y - z|}\dd z \\
&\leq 2C_{1}C_{2}2^{-j + \delta(k + 1 - j)}\log 2.
\end{align*}

(d) 当 $2^{j} - 2^{k+1} \leq |x - y| < 2^{j} + 2^{k+1}$ 时,若 $2^{k} \leq |y - z| < 2^{k+1}$,则 $|x - z| \geq |x - y| - |y - z| > 2^{j} - 2^{k+1} - 2^{k+1} \geq 2^{j} - 2^{j-1} = 2^{j-1}$.因此与 (b) 的情形一样,
\begin{equation*}
|S_{j,k}(x, y)| \leq 4C_{1}^{2}2^{-j}\log 2.
\end{equation*}

(e) 当 $|x - y| \leq 2^{j} - 2^{k+1}$ 时,若 $2^{j} \leq |x - z| < 2^{j+1}$ 且 $2^{k} \leq |y - z| < 2^{k+1}$,则 $|x - y| \geq |x - z| - |y - z| > 2^{j} - 2^{k+1}$,故 $K_{j}(x, z)\overline{K_{k}(y, z)} = 0$,
\begin{equation*}
S_{j,k}(x, y) = 0.
\end{equation*}

由 (a),(b),(c),(d),(e),
\begin{align*}
& \displaystyle\int_{\mathbb{R}}|S_{j,k}(x, y)|\dd y \\
& \leq 2(2C_{1}^{2}2^{-j}\log 2)2^{k+2} + 2(2C_{1}C_{2}2^{-j+\delta(k+1-j)}\log 2)(2^{j+1} - 2^{j} + 2^{k+2}) \\
& \quad + 2(4C_{1}^{2}2^{-j}\log 2)2^{k+2} \\
& \leq (48C_{1}^{2}\log 2)2^{k-j} + (8C_{1}C_{2}\log 2)2^{\delta}2^{\delta(k-j)} \\
& \leq 8C_{1}(6C_{1} + C_{2})(\log 2)2^{\delta}2^{-\delta|j-k|}.
\end{align*}
同样,
\begin{equation*}
\int_{\mathbb{R}}|S_{j,k}(x, y)|\dd x \leq 8C_{1}(6C_{1} + C_{2})(\log 2)2^{\delta}2^{-\delta|j-k|}.
\end{equation*}
于是与证明 $T_{j}$ 的 $L^{2}$ 有界性时一样,
\begin{equation}\tag{1.10.12}
\| T_{j}T_{k}^{*} \|_{L(H)} \leq 8C_{1}(6C_{1} + C_{2})(\log 2)2^{\delta}2^{-\delta|j-k|}.
\label{eq:1.10.12}\end{equation}

(ii) $k \geq j + 3$ 的情形.在 (i) 中把 $j$ 与 $k$ 的角色互换,可知上述估计 \eqref{eq:1.10.12} 成立.

由 (i),(ii) 的证明可见,对 $T_{j}^{*}T_{k}$ 也有同样的估计.

因此,令
\begin{equation*}
a(j) = \begin{cases} 16C_{1}^{2}\log^{2}2 & (j = -2, -1, 0, 1, 2) \\ 16C_{1}(6C_{1} + C_{2})(\log 2)2^{\delta}2^{-\delta|j|} & (|j| \geq 3), \end{cases}
\end{equation*}
则
\begin{equation*}
\| T_{j}T_{k}^{*} \|_{L(H)} + \| T_{j}^{*}T_{k} \|_{L(H)} \leq a(j - k) \quad (j, k \in \mathbb{Z}),
\end{equation*}
且
\begin{equation*}
A := \sum_{j \in \mathbb{Z}}\sqrt{a(j)} < \infty.
\end{equation*}
另一方面,前段中已证对 $f \in \mathcal{S}(\mathbb{R})$,
\begin{equation*}
\lim_{N \to \infty}\sum_{j = -N}^{N}T_{j}f(x) = Tf(x),
\end{equation*}
故由 Cotlar--Knapp--Stein 引理,
\begin{equation*}
\| Tf \|_{L^{2}} \leq A\| f \|_{L^{2}}.
\end{equation*}
\end{proof}

Hilbert 变换的积分核满足上面定理的全部假设,故重新作为

\begin{corollary}\label{cor:1.10.3}
对 Hilbert 变换 $H$,
\begin{equation*}
\| Hf \|_{L^{2}} \leq C\| f \|_{L^{2}} \quad (f \in \mathcal{S}(\mathbb{R})).
\end{equation*}
\end{corollary}
